\subsection{The stratification of culture talk}
The Simmelian formulation that equates culture talk with sociability, while revealing, also obscures a fact that is relatively uncontroversial for contemporary theorists of the culture/stratification link.   Rather than simply being a formal or neutral ``form'' of association, opportunities to engage in the type of non-instrumental sociability that Simmel points to is not likely to be equitably distributed in the social structure.  The reason for this is precisely that access to practical familiarity with the content in question is itself inequitably distributed \citep{bourdieu84}. This type of ``conversation for enjoyment'' \citep[124]{collins75}, \emph{precisely} because it is dissociated from explicit instrumental pursuits and because it is not driven by ``extrinsic rewards'' but by friendship and personal liking, comes to serve as one of the primary ways in which vertical and horizontal status-group boundaries are dynamically built and re-built.  These boundaries are recurrently recreated mainly via processes of inclusion and exclusion at the interactional level. 

Accordingly, when the very knowledge and associated experiences that are required for engaging in this form of interaction are unequally distributed \citep{reay10}---coming to form the basis for the formation of partially segregated status cultures  \citep{bourdieu84}---then we should find that members of high-status groups should be more likely to routinely engage in culture talk.  One obvious reason for this comes from what is the basic finding of the sociology of taste:  there continue to exist strong socio-demographic differences in rates of familiarity and engagement with culture and the arts, especially across groups defined by such status markers as education (and to a lesser extent, income) \citep{bennett_etal08}.  If aesthetic experiences are translated into culture talk in mundane conversational settings, then we should expect that individuals who have a higher-baseline probability of of engagement and general familiarity with culture and the arts should also be more likely to engage in culture talk.  A strong version of this argument would predict that status-linked differences in the likelihood of engaging in culture talk are purely a product of differences in cultural engagement. Under this model, status-linked differences in cultural engagement are the primary source of individual differences in cultural knowledge.  Therefore, there should not exist any residual differences across status groups \emph{within} levels of consumption.  In statistical terms, the direct effect of socio-economic status on culture talk is fully mediated by consumption frequency.

However, several considerations lead us to expect that this last version of the argument is not on the right track.  First, persons endowed with large stocks of cultural capital should be able to participate in conversation featuring culture as a topic even when they are not active consumers in the domain in question. That is, they may demonstrate what \citet[252]{peterson92} has referred as ``passing knowledge'' of a given cultural domain on the spot.  This is in fact a form of ``social skill'' that is almost a requirement among members of the young, highly-educated, new middle class \citep{dimaggio96}.   

Consistent with this hypothesis, current research on the culture-stratification link suggests that the ability to deploy cultural knowledge in interaction in partial dissociation from direct consumption is generated via multiple routes:  (1) past experience in the educational system---``scholastic capital'' in Bourdieu's \citeyearpar[59]{bourdieu84} sense---where traditional high-status culture (and increasingly ``popular'' culture) have become established part of the curriculum, (2) past experience in the upper-middle and middle-class household---``domestic capital'' in Bourdieu's \citeyearpar[59]{bourdieu84} sense---where extra-curricular forms of aesthetic engagement and arts-training have become an institutionalized element of the socialization process; (3)  routine interaction with members of the personal network who are themselves avid consumers of the cultural forms in question and who thus make culture talk a recurrent conversational topic allowing ego to gather cultural knowledge that may be used in other conversational interactions \citep{cardon_granjon05}; and finally, (4) person-to-organization connections that allow high-status individuals to participate in elite urban social circles where culture talk is a prescribed conversational resource \citep{ostrower98}.

This alternative conceptualization of the link between culture talk and social position suggests the existence of (relatively) strong structuration of culture talk by socio-economic status and other markers of difference, even after adjusting for the relative propensity to directly participate in the cultural realms in question.  These residual socio-demographic differences in culture talk propensities should be strongest precisely across those social markers most clearly associated with status and privilege, such as education and income stratum.\footnote{As noted by \citet{bennett_etal08}, this does not imply that those persons who are disengaged from cultural consumption do not have active social lives. My argument \emph{does} imply however, that the social lives of the aesthetically disengaged will be structured in distinct ways, by for instance being relatively more tied to neighborhood and locality \cite{prieur_savage11}.}   However, insofar as mechanisms associated with network connectivity and social capital may also be involved in generating the capacity to participate in ``interaction rituals''  featuring culture talk (independent of consumption), other axes of difference partially orthogonal to socio-economic status---but correlated with social connectivity---should retain their predictive power even after adjusting for levels of cultural consumption \citep{collins04}. 

Nevertheless, we should not expect all dimensions of socio-economic status to be equally good predictors of culture talk. As argued above, culture talk is deeply tied to a person's capacity to appropriate cultural goods and then mobilize discursive resources related to those consumption experiences in interaction. In this respect, culture talk is more closely connected to a given person's stock of (embodied) \emph{cultural capital} than it is to economic resources. Accordingly, we should find that those dimensions of socio-economic status that more closely track a history of accumulation of cultural capital (such as educational qualifications) should be better predictors of culture talk than those dimensions of socio-economic status that index material advantage (such as income).

\begin{quote}
Hypothesis 1:  \emph{A person's socio-economic status predicts culture-talk  frequency even after adjusting for levels of cultural participation.}
\end{quote}

Insofar as being an active participant in the arts is a clear index of cultural capital, a variation of this argument would propose that the differences in cultural talk propensities that are traceable to differences in socio-economic status themselves not independent from a person's status as an active arts participant. Instead we should expect that the predictive power of socio-economic status should change depending on whether we are talking about active or inactive culture-consumers. More specifically we should expect that status-linked differences in the likelihood of engaging in culture talk should be larger among active culture-consumers than among inactive ones.  This would be consistent with a world in which even when exposed to a similar set of experiences, persons would not only differ in the capacity to translate those experiences into a conversational resource based on social location, but this differential capacity should be most evident among persons who are more likely to engage in culture consumption in the first place.  This is consistent with a \emph{cumulative advantage} mechanism \citep{diprete_eirich06} linking consumption, status position and culture talk.  If this is correct, members of the high status groups should not only be more likely to engage in culture talk than their lower-status counterparts but this difference should be largest among the most active consumers.

\begin{quote}
Hypothesis 2:  \emph{Socio-economic status is a better predictor of culture talk frequency for culturally active persons.}
\end{quote}

\subsection{Cross-domain differences}
Does the stratification of culture talk vary across domains?  The fact that the traditional fine arts dissemination institutions, such as galleries and museums, symphony orchestras and theaters, continue to be patronized mainly by a highly selected, educated elite is a well-established empirical finding \citep{dimaggio_mukhtar04}.  We also know that attempts to ``democratize'' access to the fine arts on the part of these institutions have met with equivocal results, and in most cases have failed outright \citep{alexander96}.  The dependence of non-profit arts organizations on this consumer base provides high-status elites with recurrent occasions to meet others in the same circles at sponsored social gatherings \citep{ostrower98}.  These are interaction opportunities that are ideally suited for the deployment of culture talk.  

In addition, it is clear that the field of fine arts production possess a long-standing legacy of discursive tools---developed by professional critics and lay enthusiasts---for the critical evaluation and the \emph{intellectualized} expression of taste judgments \citep{weber77}.  The availability of established linguistic conventions with which to categorize and evaluate cultural products may facilitate the mobilization of this type of content as a conversational resource among high-status elites \citep{dimaggio87}.  In all, if we presume that routine access to and knowledge about the visual and performance arts is certainly more strongly stratified by education and socio-economic status than is the case for more commercialized domains, then we should find that culture talk involving this type of content should also be more highly stratified by status.

\begin{quote}
Hypothesis 3: \emph{Socioeconomic status is a better predictor of culture-talk frequency in fine arts domains in comparison to commercial arts domains.}
\end{quote}

Recent studies of the consumption of the popular arts, inspired by Bourdieu's \citeyearpar{bourdieu93a} work on ``fields of cultural production,'' has begun to conceive of the popular and commercial arts as equally capable of being the subject of attempts to impose hierarchies of value and taste, the construction of exclusionary canons, and the establishment systems of cultural currency as the fine arts. Rather than operating according to a logic that is antithetical to that which structures fine arts consumption fields (as presumed by an earlier generation of popular culture studies, the realm of commercial culture production and consumption comes to be structured according to hierarchies of perception and evaluation that are premised upon categories of perception and vocabularies of appreciation developed in autonomous fields of cultural production to the products of the culture industry.  In this respect, ``artistic hierarchies, which rank producers according to their aesthetic or expressive value{\ldots}are becoming a central structuring force in a growing number of fields of production'' including that of rock and roll, a quintessential instance of the commercial arts \citep[98]{regev94}.  

More recently, other authors have begun to extend this analysis of the dynamics whereby other popular culture fields-such as film---come to acquire a hierarchical form modeled on the artistic field \citep[e.g.][]{allen_lincoln04}. Thus, as \citet{baumann01} has argued, certain commercial forms (such as film) have developed an intellectualized vocabulary of appreciation that matches that of their fine arts counterparts. The adoption of this ideology of autonomous art by producers and gate-keepers of the popular arts mirrors the increasing propensity of  (culturally advantaged) consumers of popular culture to adopt a similar stance \citep{thornton96}. 

Just like in the fine arts field, where struggles to define artistic value based on what is defined as `art' and `not art' are constantly in play, ``a similar use of accumulated knowledge and discriminatory skill is apparent in low cultural forms, and has the same hierarchical effect.  Low culture, that is to say, generates its own capital---most obviously perhaps, in those forms{\ldots}which are organized around  exclusiveness, but equally significant for the fans:{\ldots}of even the most inclusive forms'' \citep[9]{frith98}. Insofar as the popular arts are subject to this intellectualization process, and insofar as critical gatekeepers develop a language of appreciation and historical canonization, then they will be likely to figure as fodders for culture talk among high-status classes. Thus, in the very same way in which both the commercial and the traditional arts can be appropriated ``aesthetically'' by members of high status classes \citep{peterson_kern96}, we should find that high status groups are equal opportunity culture talkers across domains:

\begin{quote}
Hypothesis 4:  \emph{Socioeconomic status predicts culture-talk frequency across all cultural domains.}
\end{quote}

If culture talk is spurred by the extent to which a person is familiar with intellectualized discourses of cultural appreciation, then we should also expect that there should exist predictable \emph{asymmetries} in the extent to which the consumption of fine arts spurs cultural talk in commercial domains as opposed to the reverse. More specifically, we should find that the consumption of the fine arts (as an index of familiarity with the discursive conventions of art) should be an equally good predictor of culture talk in so-called commercial domains as the consumption of commercial art, but not vice versa. That is, fine arts consumers should be as (or more) likely to engage in culture talk in commercial domains, but commercial arts consumers should be less likely than fine arts consumers to engage in culture talk in fine arts domains:

\begin{quote}
Hypothesis 5:  \emph{Fine arts consumption is a better predictor of culture talk across both commercial and fine arts domains than commercial arts consumption.}
\end{quote}

